\documentclass[11pt,a4paper]{article}

\usepackage[utf8]{inputenc}
\usepackage[T1]{fontenc}
\usepackage[french]{babel}
\usepackage[a4paper,margin=2.2cm]{geometry}
\usepackage{amsmath,amssymb,amsthm,mathtools}
\usepackage{lmodern}
\usepackage{enumitem}
\usepackage{hyperref}
\usepackage{xcolor}
\usepackage{fancyhdr}
\usepackage{tikz}

\hypersetup{
    colorlinks=true,
    linkcolor=black,
    urlcolor=blue
}

\setlength{\parindent}{0pt}
\setlength{\parskip}{0.6em}

\pagestyle{fancy}
\fancyhf{}
\fancyhead[L]{Fiche d’exercices --- Algèbre linéaire Chapitre 2}
\fancyhead[R]{Pratique calculatoire}
\fancyfoot[C]{\thepage}

\begin{document}

\begin{center}
    {\LARGE \textbf{Fiche d’exercices --- Chapitre 2}}\\[0.4em]
    {\Large \textbf{Pratique calculatoire et objets algébriques de base}}\\[0.6em]
    Algèbre linéaire
\end{center}

\textbf{AIDE de lecture type d’exo :} APP = application directe, CALC = calcul rédigé et méthode, DEM = démonstration, PREP = réflexion plus poussée. Les étoiles représentent un gradient de difficulté de 1 à 5.

\tableofcontents

\newpage

\section*{Exercices d’application directe}
\addcontentsline{toc}{section}{Exercices d’application directe}

\textbf{Exercice 1 [APP $\star$]}\\
\textbf{Réduction d’expressions}\\
Réduire les expressions suivantes :
\begin{enumerate}[label=\alph*)]
    \item $2x+3x-4$,
    \item $7a-2a+5$,
    \item $3x-2+4x+9$,
    \item $5y-(2y-3)$,
    \item $-3(2x-1)+4x$.
\end{enumerate}

\textbf{Exercice 2 [APP $\star$]}\\
\textbf{Développement simple}\\
Développer et réduire :
\begin{enumerate}[label=\alph*)]
    \item $3(x+2)$,
    \item $-2(x-5)$,
    \item $(x+1)(x+4)$,
    \item $(2x-3)(x+5)$,
    \item $(a-b)(a+2b)$.
\end{enumerate}

\textbf{Exercice 3 [APP $\star$]}\\
\textbf{Identités remarquables immédiates}\\
Développer :
\begin{enumerate}[label=\alph*)]
    \item $(x+3)^2$,
    \item $(2x-1)^2$,
    \item $(a-b)^2$,
    \item $(x+5)(x-5)$,
    \item $(3a+b)^2$.
\end{enumerate}

\textbf{Exercice 4 [APP $\star\star$]}\\
\textbf{Factorisation élémentaire}\\
Factoriser :
\begin{enumerate}[label=\alph*)]
    \item $4x+4$,
    \item $3x^2-6x$,
    \item $x^2-9$,
    \item $a^2-2ab+b^2$,
    \item $5x^2+10x$.
\end{enumerate}

\textbf{Exercice 5 [APP $\star\star$]}\\
\textbf{Fractions numériques}\\
Calculer et simplifier :
\begin{enumerate}[label=\alph*)]
    \item $\dfrac{2}{3}+\dfrac{1}{6}$,
    \item $\dfrac{5}{4}-\dfrac{3}{8}$,
    \item $\dfrac{2}{5}\times\dfrac{15}{4}$,
    \item $\dfrac{7}{9}\div\dfrac{14}{3}$,
    \item $\dfrac{3}{2}+\dfrac{5}{3}-1$.
\end{enumerate}

\textbf{Exercice 6 [APP $\star\star$]}\\
\textbf{Fractions littérales et conditions d’existence}\\
Simplifier, en précisant les conditions d’existence :
\begin{enumerate}[label=\alph*)]
    \item $\dfrac{x^2-1}{x-1}$,
    \item $\dfrac{x^2-4}{x+2}$,
    \item $\dfrac{1}{x}+\dfrac{1}{x+1}$,
    \item $\dfrac{2x}{x^2+x}$,
    \item $\dfrac{x^2-9}{x^2-3x}$.
\end{enumerate}

\textbf{Exercice 7 [APP $\star\star$]}\\
\textbf{Puissances}\\
Simplifier en utilisant les règles usuelles :
\begin{enumerate}[label=\alph*)]
    \item $x^2x^5$,
    \item $\dfrac{x^7}{x^3}$,
    \item $(x^2)^4$,
    \item (2x)^3,
    \item $a^{-2}a^5$.
\end{enumerate}

\textbf{Exercice 8 [APP $\star\star$]}\\
\textbf{Valeur absolue}\\
Calculer :
\begin{enumerate}[label=\alph*)]
    \item $|5|$,
    \item $|-7|$,
    \item $|3-8|$,
    \item $|-2+6|$,
    \item $|x|$ lorsque $x=-4$.
\end{enumerate}

\textbf{Exercice 9 [APP $\star\star$]}\\
\textbf{Équations produits}\\
Résoudre dans $\mathbb{R}$ :
\begin{enumerate}[label=\alph*)]
    \item $(x-2)(x+3)=0$,
    \item $(2x+1)(x-4)=0$,
    \item $x^2-16=0$,
    \item $(x+5)^2=0$,
    \item $(3x-6)(x+1)=0$.
\end{enumerate}

\textbf{Exercice 10 [APP $\star\star$]}\\
\textbf{Factorisations par regroupement}\\
Factoriser :
\begin{enumerate}[label=\alph*)]
    \item $ax+ay+bx+by$,
    \item $x^2+3x+2x+6$,
    \item $a(x-1)+b(x-1)$,
    \item $2x^2-2x+3x-3$,
    \item $(u+v)a+(u+v)b$.
\end{enumerate}

\newpage

\section*{Exercices de calcul rédigé et de méthode}
\addcontentsline{toc}{section}{Exercices de calcul rédigé et de méthode}

\textbf{Exercice 11 [CALC $\star\star$]}\\
\textbf{Identité ou équation ?}\\
Pour chacune des affirmations suivantes, dire s’il s’agit d’une identité ou d’une équation :
\begin{enumerate}[label=\alph*)]
    \item $(x+1)^2=x^2+2x+1$,
    \item $x^2=9$,
    \item $(a-b)^2=a^2-2ab+b^2$,
    \item $2x+3=7$,
    \item $(x+y)(x-y)=x^2-y^2$.
\end{enumerate}

\textbf{Exercice 12 [CALC $\star\star$]}\\
\textbf{Choisir entre développer et factoriser}\\
Pour chacune des expressions suivantes, indiquer quelle méthode semble la plus pertinente selon l’objectif demandé, puis effectuer le calcul :
\begin{enumerate}[label=\alph*)]
    \item simplifier $(x+2)(x-3)+(x+2)(x+5)$ ;
    \item résoudre $x^2-25=0$ ;
    \item simplifier $(x-1)^2-(x+1)^2$ ;
    \item simplifier $\dfrac{x^2-4}{x-2}$ ;
    \item résoudre $(x+3)^2-(x+3)=0$.
\end{enumerate}

\textbf{Exercice 13 [CALC $\star\star\star$]}\\
\textbf{Attention aux parenthèses}\\
Développer et réduire soigneusement :
\begin{enumerate}[label=\alph*)]
    \item $-(x-3)+2(x+1)$,
    \item $-(2x-5)-(x+4)$,
    \item $(3-x)(x+2)$,
    \item $(1-2x)^2$,
    \item $-(x-1)(x+1)$.
\end{enumerate}

\textbf{Exercice 14 [CALC $\star\star\star$]}\\
\textbf{Expressions rationnelles}\\
Simplifier, en détaillant chaque étape et en donnant l’ensemble de définition :
\begin{enumerate}[label=\alph*)]
    \item $\dfrac{1}{x}+\dfrac{2}{x+1}$,
    \item $\dfrac{x+1}{x}-\dfrac{1}{x}$,
    \item $\dfrac{x^2-1}{x^2+x}$,
    \item $\dfrac{1}{x-1}-\dfrac{1}{x+1}$,
    \item $\dfrac{x^2-4}{x^2-2x}$.
\end{enumerate}

\textbf{Exercice 15 [CALC $\star\star\star$]}\\
\textbf{Substitution utile}\\
En posant $u=x+1$, simplifier ou résoudre :
\begin{enumerate}[label=\alph*)]
    \item $(x+1)^2+3(x+1)+2$,
    \item $(x+1)^2-4$,
    \item $(x+1)^2-5(x+1)+6=0$,
    \item $(x+1)^2-(x+1)$,
    \item $\dfrac{1}{x+1}+\dfrac{2}{(x+1)^2}$.
\end{enumerate}

\textbf{Exercice 16 [CALC $\star\star\star$]}\\
\textbf{Carrés et différences de carrés}\\
Simplifier sans développer complètement lorsque cela est pertinent :
\begin{enumerate}[label=\alph*)]
    \item $(x+4)^2-(x-4)^2$,
    \item $(2x+1)^2-(2x-1)^2$,
    \item $(a+b)^2-(a-b)^2$,
    \item $(x+3)(x-3)+(x-3)^2$,
    \item $(x+1)^2-x^2$.
\end{enumerate}

\textbf{Exercice 17 [CALC $\star\star\star$]}\\
\textbf{Valeur absolue et écriture par cas}\\
Écrire sans valeur absolue selon les cas :
\begin{enumerate}[label=\alph*)]
    \item $|x-2|$,
    \item $|3-x|$,
    \item $|2x+1|$,
    \item $|x|+|x-1|$,
    \item $|x+2|-|x-2|$.
\end{enumerate}

\textbf{Exercice 18 [CALC $\star\star\star$]}\\
\textbf{Premiers polynômes}\\
Pour chacun des polynômes suivants, donner le degré, puis chercher d’éventuelles racines évidentes :
\begin{enumerate}[label=\alph*)]
    \item $P(x)=3x^2-12x$,
    \item $Q(x)=x^2-5x+6$,
    \item $R(x)=x^3-x$,
    \item $S(x)=x^2-1$,
    \item $T(x)=x^3-3x^2+2x$.
\end{enumerate}

\textbf{Exercice 19 [CALC $\star\star\star$]}\\
\textbf{Équations algébriques simples}\\
Résoudre dans $\mathbb{R}$ :
\begin{enumerate}[label=\alph*)]
    \item $x^2-5x+6=0$,
    \item $x^2-x=0$,
    \item $x^3-4x=0$,
    \item $(x-1)^2-(x-1)=0$,
    \item $(x+2)^2-9=0$.
\end{enumerate}

\textbf{Exercice 20 [CALC $\star\star\star$]}\\
\textbf{Binôme de Newton, petits cas}\\
Développer :
\begin{enumerate}[label=\alph*)]
    \item $(a+b)^3$,
    \item $(x-1)^3$,
    \item $(2x+3)^3$,
    \item $(u-v)^3$,
    \item $(x+2)^4$.
\end{enumerate}

\newpage

\section*{Exercices de démonstration}
\addcontentsline{toc}{section}{Exercices de démonstration}

\textbf{Exercice 21 [DEM $\star\star$]}\\
\textbf{Simplification légitime}\\
Montrer que si $ac=bc$ avec $c\neq 0$, alors $a=b$.

\textbf{Exercice 22 [DEM $\star\star$]}\\
\textbf{Produit nul}\\
Montrer que pour tous réels $a$ et $b$,
\[
ab=0 \Longleftrightarrow a=0 \text{ ou } b=0.
\]

\textbf{Exercice 23 [DEM $\star\star$]}\\
\textbf{Différence de deux carrés}\\
Montrer que pour tous réels $a$ et $b$,
\[
a^2-b^2=(a-b)(a+b).
\]

\textbf{Exercice 24 [DEM $\star\star$]}\\
\textbf{Mise en facteur commun}\\
Montrer que pour tous réels $a,b,c$,
\[
ab+ac=a(b+c).
\]

\textbf{Exercice 25 [DEM $\star\star\star$]}\\
\textbf{Somme des premiers entiers}\\
Montrer que pour tout $n\in\mathbb{N}$,
\[
1+2+\cdots+n=\frac{n(n+1)}{2}.
\]

\textbf{Exercice 26 [DEM $\star\star\star$]}\\
\textbf{Somme géométrique}\\
Montrer que pour tout réel $q\neq 1$ et tout $n\in\mathbb{N}$,
\[
1+q+q^2+\cdots+q^n=\frac{1-q^{n+1}}{1-q}.
\]

\textbf{Exercice 27 [DEM $\star\star\star$]}\\
\textbf{Valeur absolue et carré}\\
Montrer que pour tout réel $x$,
\[
x^2=|x|^2.
\]

\textbf{Exercice 28 [DEM $\star\star\star$]}\\
\textbf{Inégalité triangulaire, cas simple}\\
En utilisant le cours, montrer que pour tous réels $x,y$,
\[
|x+y|\leq |x|+|y|.
\]

\textbf{Exercice 29 [DEM $\star\star\star$]}\\
\textbf{Un calcul utile}\\
Montrer que pour tout réel $x$,
\[
(x+1)^2-(x-1)^2=4x.
\]

\textbf{Exercice 30 [DEM $\star\star\star$]}\\
\textbf{Factorisation d’un polynôme simple}\\
Montrer que pour tout réel $x$,
\[
x^3-x=x(x-1)(x+1).
\]

\textbf{Exercice 31 [DEM $\star\star\star\star$]}\\
\textbf{Somme de nombres impairs}\\
Montrer que pour tout $n\in\mathbb{N}$,
\[
1+3+5+\cdots+(2n+1)=(n+1)^2.
\]

\textbf{Exercice 32 [DEM $\star\star\star\star$]}\\
\textbf{Racine évidente et factorisation}\\
Soit
\[
P(x)=x^3-6x^2+11x-6.
\]
Vérifier que $1$, $2$ et $3$ sont racines de $P$, puis en déduire une factorisation complète de $P(x)$.

\newpage

\section*{Exercices de réflexion}
\addcontentsline{toc}{section}{Exercices de réflexion}

\textbf{Exercice 33 [PREP $\star\star\star$]}\\
\textbf{Un faux réflexe classique}\\
Dire si l’affirmation suivante est vraie ou fausse :
\[
(a+b)^2=a^2+b^2.
\]
Si elle est fausse, donner un contre-exemple puis écrire la bonne formule.

\textbf{Exercice 34 [PREP $\star\star\star$]}\\
\textbf{Simplification dangereuse}\\
Résoudre dans $\mathbb{R}$ l’équation
\[
x(x-1)=x(x+2),
\]
sans jamais simplifier abusivement.

\textbf{Exercice 35 [PREP $\star\star\star$]}\\
\textbf{Écriture astucieuse}\\
Simplifier :
\[
\frac{1}{x}-\frac{1}{x+1}
\]
puis déterminer le signe du résultat lorsque $x>0$.

\textbf{Exercice 36 [PREP $\star\star\star$]}\\
\textbf{Factoriser ou développer ?}\\
On considère
\[
A=(x-2)(x+5)-(x-2)(x-1).
\]
\begin{enumerate}[label=\alph*)]
    \item Simplifier $A$ de la manière la plus intelligente possible.
    \item Retrouver le même résultat en développant tout.
\end{enumerate}

\textbf{Exercice 37 [PREP $\star\star\star\star$]}\\
\textbf{Expression symétrique}\\
Simplifier au maximum :
\[
(x+y)^2+(x-y)^2.
\]
Puis proposer une expression analogue pour
\[
(x+y)^2-(x-y)^2.
\]

\textbf{Exercice 38 [PREP $\star\star\star\star$]}\\
\textbf{Équation avec paramètre caché}\\
Résoudre dans $\mathbb{R}$ :
\[
(x+1)^2-3(x+1)+2=0.
\]
On demandera de faire apparaître clairement un changement d’écriture adapté.

\textbf{Exercice 39 [PREP $\star\star\star\star$]}\\
\textbf{Somme et produit connus}\\
Déterminer deux réels $u$ et $v$ tels que
\[
u+v=7 \qquad \text{et} \qquad uv=10.
\]
Mettre le problème sous la forme d’une équation algébrique simple.

\textbf{Exercice 40 [PREP $\star\star\star\star$]}\\
\textbf{Valeur absolue et distance}\\
Résoudre dans $\mathbb{R}$ :
\[
|x-3|=5.
\]
Puis interpréter géométriquement le résultat sur la droite réelle.

\textbf{Exercice 41 [PREP $\star\star\star\star$]}\\
\textbf{Somme géométrique concrète}\\
Calculer
\[
1+2+2^2+\cdots+2^{10}
\]
sans effectuer les additions une à une.

\textbf{Exercice 42 [PREP $\star\star\star\star\star$]}\\
\textbf{Un calcul qui télescope presque}\\
Simplifier, pour $x\neq 0$ et $x\neq -1$,
\[
\frac{1}{x}-\frac{1}{x+1}.
\]
Puis en déduire une méthode pour calculer rapidement :
\[
\left(1-\frac12\right)+\left(\frac12-\frac13\right)+\left(\frac13-\frac14\right)+\cdots+\left(\frac1n-\frac1{n+1}\right).
\]

\textbf{Exercice 43 [PREP $\star\star\star\star\star$]}\\
\textbf{Une factorisation moins immédiate}\\
Factoriser
\[
x^4-1.
\]
On demandera une factorisation complète sur $\mathbb{R}$.

\textbf{Exercice 44 [PREP $\star\star\star\star\star$]}\\
\textbf{Polynôme et racines évidentes}\\
On considère
\[
P(x)=x^3+x^2-4x-4.
\]
\begin{enumerate}[label=\alph*)]
    \item Montrer que $-1$ est une racine de $P$.
    \item En déduire une factorisation de $P(x)$.
    \item Résoudre $P(x)=0$ dans $\mathbb{R}$.
\end{enumerate}

\textbf{Exercice 45 [PREP $\star\star\star\star\star$]}\\
\textbf{Problème de synthèse final}\\
Pour chacune des affirmations suivantes :
\begin{enumerate}[label=\alph*)]
    \item dire si elle est vraie ou fausse ;
    \item si elle est vraie, la démontrer ;
    \item si elle est fausse, donner un contre-exemple.
\end{enumerate}

\begin{enumerate}[label=\arabic*)]
    \item Si $a+b=0$, alors $a^2=b^2$.
    \item Si $a^2=b^2$, alors $a=b$.
    \item Si $x^2-5x+6=0$, alors $x=2$ ou $x=3$.
    \item Si $\dfrac{ab}{ac}=\dfrac{b}{c}$, alors cette simplification est toujours licite.
\end{enumerate}

\newpage

\section*{Pour aller plus loin}
\addcontentsline{toc}{section}{Pour aller plus loin}

\textbf{Exercice 46 [PREP $\star\star\star\star\star$]}\\
\textbf{Expression à paramètre}\\
Pour $a\in\mathbb{R}$, résoudre dans $\mathbb{R}$ :
\[
(x-a)^2-(x-a)=0.
\]

\textbf{Exercice 47 [PREP $\star\star\star\star\star$]}\\
\textbf{Reconnaître une structure}\\
Simplifier :
\[
(x+2)^2-4(x+2)+4.
\]
Que remarque-t-on ?

\textbf{Exercice 48 [PREP $\star\star\star\star\star$]}\\
\textbf{Somme de carrés et positivité}\\
Montrer que pour tout réel $x$,
\[
x^2+2x+2>0.
\]
On demandera une rédaction mettant en évidence une écriture adaptée.

\textbf{Exercice 49 [DEM $\star\star\star\star$]}\\
\textbf{Somme des carrés}\\
Admettre la formule
\[
1^2+2^2+\cdots+n^2=\frac{n(n+1)(2n+1)}{6},
\]
puis la vérifier pour $n=1$, $n=2$ et $n=3$.\\
Comparer ensuite son ordre de grandeur avec la somme
\[
1+2+\cdots+n.
\]

\textbf{Exercice 50 [PREP $\star\star\star\star\star$]}\\
\textbf{Synthèse générale de calcul}\\
Simplifier au maximum :
\[
A=\frac{x^2-1}{x-1}-\frac{x^2-4}{x-2},
\]
en précisant soigneusement les conditions d’existence.\\
Puis résoudre l’équation $A=0$.

\end{document}